

\begin{document}
	
	% RESET all counters at the beginning
	\setcounter{section}{0}
	\setcounter{subsection}{0}
	\setcounter{subsubsection}{0}
	\setcounter{paragraph}{0}
	
	% Depth for numbering and TOC
	\setcounter{secnumdepth}{1}  % Number sections only
	% Part in TOC: footnotesize, bold, NO page break
	\makeatletter
	\renewcommand*\l@part[2]{%
		\ifnum \c@tocdepth >-2\relax
		\addpenalty{-\@highpenalty}%
		\addvspace{0.8em \@plus\p@}%
		{\leftskip 0em \relax
			\rightskip \@tocrmarg
			\parfillskip -\rightskip
			\parindent 0em \relax\@afterindenttrue
			\interlinepenalty\@M
			\leavevmode
			{\footnotesize\bfseries #1}\nobreak
			\leaders\hbox{$\m@th\mkern \@dotsep mu\hbox{}\mkern \@dotsep mu$}\hfill
			\nobreak\hb@xt@\@pnumwidth{\hss #2}\par}%
		\addvspace{0.2em \@plus\p@}%
		\nobreak
		\fi}
	\makeatother
	
	% Chapter in TOC: footnotesize, bold
	\renewcommand{\cftchapfont}{\footnotesize\bfseries}
	\renewcommand{\cftchappagefont}{\footnotesize\bfseries}
	\setlength{\cftbeforechapskip}{0.3em}
	
	% Only Chapters in TOC (no Sections/Subsections)
	\setcounter{tocdepth}{0}
	

	\begin{titlepage}
		\centering
		\vspace*{2cm}
	{\Huge\textbf{Fundamentale Fraktal-Geometrische Feldtheorie (FFGFT) oder T0-Theorie: Zeit-Masse-Dualität}}\\[1cm]
	{\Large Teil 7: Gitter im Hilbertraum, Spektraltheorie und Galois-Struktur}\\[2cm]
		{\large Johann Pascher}\\[1cm]
		{\large 2026}
		\vfill
	\end{titlepage}

	\frontmatter
	\pagestyle{fancy}
	\makeatletter
	\let\ps@plain\ps@fancy
	\let\ps@empty\ps@fancy
	\makeatother

	\mainmatter
	\pagestyle{fancy}
	\tableofcontents

	\chapter{Einleitung zu Band 7}
	% =============================================================================
% EINLEITUNG ZU BAND 7: GITTER IM HILBERTRAUM, SPEKTRALTHEORIE UND
% GALOIS-STRUKTUR (Bd. 7 von 8)
% =============================================================================

\section*{Siebter Band der Dokumentensammlung}

Band~7 setzt die in Band~6 begonnene Fortführung der Sammlung fort und umfasst
die Dokumente 311 bis 343 aus dem August und dem Beginn des Septembers 2026.
In diesem Zeitraum verlagert sich der Schwerpunkt des Korpus von der
Projektion in den Beobachtungsraum auf die algebraische Struktur der kompakten
Geometrie selbst.

\subsection*{Inhalt und Schwerpunkt}

Der erste Teil beginnt mit der Frage, wie ein vierdimensionales Gitter auf drei
Raumdimensionen kommt (Dok.~311), und führt über die Abschlussskala, in der
$\Lambda$ als abgeleitete Größe der kompakten Geometrie erscheint, und die
thermische Geschichte auf dem Zeitzyklus zum Gitter im Hilbertraum
(Dok.~314). Dort wird das Gitter als Indexmenge der Fouriermoden gelesen;
Geometrie wird zu Spektrum und Faser. Die Untersuchung der Form von
$K_\text{frak}$ und der Vergleich mit der Riemannschen Zeta-Funktion schließen
diesen Teil ab.

Der zweite Teil bündelt die Leptonen-Generationen in topologischer Lesart, den
Geltungsbereich der Masseableitungen, die Spektraltheorie der FFGFT-Tori mit
ihrer Hilbertraum-Abbildung, die algebraische Herleitung der
$SU(3)_c$-Eichstruktur aus der $\mathbb{Z}_3$-Trialität und die Herleitung des
Weinberg-Winkels bei $M_Z$ mit vollständigem Renormierungsgruppenlauf. Im
dritten Teil stehen der Vergleich mit der Hyperbit-Vakuumstruktur
$G(6,\mathbb{Z}_3\mathbb{C})$, der FFGFT-Hawking-Mechanismus, die
Selbstadjungiertheit des fundamentalen Operators $\hat F$, die
Kopplungsregime von Resonanz und Synchronisation sowie die zwei Wege vom
Diskreten zum Kontinuum.

Der vierte Teil enthält die Galois-Brücke. Ausgehend vom endlichen Körper
$GF(9)$ (Dok.~336) wird gezeigt, dass Frobenius-Automorphismus und
FFGFT-Sektorpaarung strukturgleich sind; über $GF(27)^*$ folgen die
Gegenüberstellung der Massenverhältnisse mit Galois-Gruppenordnungen, die
Frobenius-Trennung von massivem und masselosem Sektor, die
Neutrino-Massenhierarchie, die Faktorisierungsklassen in $G(6)$ mit der
Harmonik der Primzahlen und die Zeta-Funktion des $T^4/\mathbb{Z}_3$-Torus.

\subsection*{Gliederung}

Die 33 Dokumente sind in vier Teile gegliedert: Vier auf drei,
Abschlussskala und Gitter im Hilbertraum (Dok.~311--316); Leptonen,
Spektraltheorie und Eichstruktur (Dok.~317--323); Hyperbits,
Hawking-Mechanismus, Kopplung und Kontinuum (Dok.~324--335); die
Galois-Brücke (Dok.~336--343).

\subsection*{Zum Lesen der Dokumente}

Die Dokumente 311 bis 314 sind ursprünglich als mehrkapitelige Abhandlungen
gesetzt; im Buch erscheinen ihre Kapitel als Abschnitte des jeweiligen
Dokuments. Statusmarker und Korrekturregister gelten wie in den
vorangehenden Bänden; maßgeblich ist jeweils die jüngste Fassung eines
Ergebnisses im Korpus.

\subsection*{Stand}

Die Dokumente dieses Bandes tragen den Stand ihrer jeweiligen Fassung im
Korpus zum Oktober 2026.


	% --- part ---
	\part{Teil I --- Vier auf drei, Abschlussskala und Gitter im Hilbertraum}
	\FFGFTdocDemote{Dok.~311 --- Vier auf drei --- Wie kommt ein 4D-Gitter auf drei Raumdimensionen?}{../ch/311_Vier_auf_Drei_De_ch}
	\FFGFTdocDemote{Dok.~312 --- Die Abschlussskala: $\Lambda$ als abgeleitete Größe der kompakten Geometrie}{../ch/312_Abschlussskala_De_ch}
	\FFGFTdoc{Anhang A zu Dok.~312 --- Altermagnetismus als Laboranalogie}{../ch/312_Anhang_Altermagnet_De_ch}
	\FFGFTdocDemote{Dok.~313 --- Kein Anfang: Die thermische Geschichte auf dem Zeitzyklus}{../ch/313_Kein_Anfang_De_ch}
	\FFGFTdocDemote{Dok.~314 --- Das Gitter im Hilbertraum: Geometrie, Spektrum und Faser}{../ch/314_Gitter_im_Hilbertraum_De_ch}
	\FFGFTdoc{Dok.~315 --- Die Form von $K_\text{frak}$: additiv oder multiplikativ?}{../ch/315_Kfrak_Form_De_ch}
	\FFGFTdoc{Dok.~316 --- FFGFT und die Riemannsche Zeta-Funktion}{../ch/316_Riemann_Zeta_De_ch}
	% --- part ---
	\part{Teil II --- Leptonen, Spektraltheorie und Eichstruktur}
	\FFGFTdoc{Dok.~317 --- Topologischer Ursprung der Leptonen-Generationen: KSAU und FFGFT}{../ch/317_KSAU_FFGFT_Leptonen_De_ch}
	\FFGFTdoc{Dok.~318 --- Geltungsbereich der Masseableitungen in FFGFT: Ruhemassen, Hadrondynamik und offene Brücken}{../ch/318_Masseableitung_Geltungsbereich_De_ch}
	\FFGFTdoc{Dok.~319 --- Das Proton als schwingender Torus: Geometrische Grundlagen der Massenemergenz in der FFGFT}{../ch/319_Proton_Torus_De_ch}
	\FFGFTdoc{Dok.~320 --- Ausführliche Feldgeometrische Spektraltheorie}{../ch/320_Spektraltheorie_De_ch}
	\FFGFTdoc{Dok.~321 --- Algebraische Herleitung der $SU(3)_c$-Eichstruktur}{../ch/321_SU3_Z3_Emergenz_De_ch}
	\FFGFTdoc{Dok.~322 --- Spektraltheorie und Hilbertraum-Abbildung der FFGFT-Tori}{../ch/322_Spektraltheorie_Hilbert_De_ch}
	\FFGFTdoc{Dok.~323 --- Herleitung des Weinberg-Winkels bei $M_Z$ aus der FFGFT-Geometrie}{../ch/323_Weinberg_Winkel_RGE_De_ch}
	% --- part ---
	\part{Teil III --- Hyperbits, Hawking-Mechanismus, Kopplung und Kontinuum}
	\FFGFTdoc{Dok.~324 --- G(6,Z$_3$C)-Vakuumstruktur und FFGFT}{../ch/324_G6_Z3C_FFGFT_Vergleich_De_ch}
	\FFGFTdoc{Dok.~325 --- Der FFGFT-Hawking-Mechanismus}{../ch/325_Hawking_FFGFT_De_ch}
	\FFGFTdoc{Dok.~326 --- Schwarze Löcher aus Hyperbits und aus FFGFT}{../ch/326_Matzke_FFGFT_Vergleich_De_ch}
	\FFGFTdoc{Dok.~327 --- Selbstadjungiertheit des fundamentalen fraktalen Operators $\hat{F}$}{../ch/327_Selbstadjungiertheit_F_De_ch}
	\FFGFTdoc{Dok.~328 --- Kopplungsregime, Resonanz und Synchronisation}{../ch/328_Kopplungsregime_Resonanz_De_ch}
	\FFGFTdoc{Dok.~329 --- Ist die Stabilitätsschwelle $n_{\mathrm{thresh}}$ universell?}{../ch/329_nthresh_Skalenanalyse_De_ch}
	\FFGFTdoc{Dok.~330 --- Drei Operationen von $T^4$ zum Beobachtraum}{../ch/330_Operationen_T4_Abgrenzung_De_ch}
	\FFGFTdoc{Dok.~332 --- Zwei Wege vom Diskreten zum Kontinuum}{../ch/332_Kontinuumskonstruktionen_De_ch}
	\FFGFTdoc{Dok.~333 --- $K_{\text{frak}}$, Ausrollen und die $\tilde{T}\cdot m$-Dualität}{../ch/333_Kfrak_Dualitaet_De_ch}
	\FFGFTdoc{Dok.~334 --- Superposition ohne Zeit --- Was die Abwesenheit des Zeitoperators in der QM bedeutet}{../ch/334_Superposition_ohne_Zeit_De_ch}
	\FFGFTdocDemote{Dok.~335 --- Massegebundene Zeitskalen und Interferometrie}{../ch/335_Frequenzskalen_De_ch}
	% --- part ---
	\part{Teil IV --- Die Galois-Brücke: GF(9), GF(27) und die Zeta-Funktion des Torus}
	\FFGFTdoc{Dok.~336 --- Die algebraische Brücke: FFGFT und $\mathbb{Z}_3\mathbb{C}$}{../ch/336_GF9_FFGFT_Bruecke_De_ch}
	\FFGFTdoc{Dok.~337 --- Zeit-Emergenz: FFGFT und GALG im Vergleich}{../ch/337_Zeit_Emergenz_FFGFT_GALG_De_ch}
	\FFGFTdoc{Dok.~338 --- Masseberechnungen und Galois-Gegenüberstellung}{../ch/338_Galois_Massen_FFGFT_De_ch}
	\FFGFTdoc{Dok.~339 --- Frobenius-Trennung massiv/masselos in $GF(27)^*$}{../ch/339_Frobenius_Trennung_De_ch}
	\FFGFTdoc{Dok.~340 --- Neutrino-Massenhierarchie aus der $GF(27)^*$-Frobenius-Struktur}{../ch/340_Neutrino_Galois_De_ch}
	\FFGFTdoc{Dok.~341 --- GF$(27)$ in GALG und FFGFT: Warum G$(3)$ kein GF$(27)$ enthält und G$(6)$ schon}{../ch/341_GF27_GALG_FFGFT_De_ch}
	\FFGFTdoc{Dok.~342 --- Faktorisierungsklassen in G$(6)$ und die Harmonik der Primzahlen}{../ch/342_Faktorisierung_Harmonik_De_ch}
	\FFGFTdoc{Dok.~343 --- Die Zeta-Funktion des T$^4$/Z$_3$-Torus: Orbifold-Projektion und Dirichlet-L-Funktionen}{../ch/343_Zeta_Galois_FFGFT_De_ch}
\end{document}
